\documentclass[11pt]{article}
\usepackage[margin=1in]{geometry}
\usepackage[T1]{fontenc}
\usepackage{lmodern,amsmath,amssymb,amsthm,mathtools,hyperref}
\setlength{\emergencystretch}{3em}
\newtheorem{theorem}{Theorem}[section]
\newtheorem{lemma}[theorem]{Lemma}
\newtheorem{corollary}[theorem]{Corollary}
\newtheorem{proposition}[theorem]{Proposition}
\newcommand{\E}{\mathbb E}
\newcommand{\tr}{\operatorname{tr}}
\title{Nonlinear seed encodings, adaptive weights,\\and entropy-order transitions}
\author{Joshua Mark\thanks{ORCID: \href{https://orcid.org/0009-0003-5905-2382}{0009-0003-5905-2382}. AI tools assisted with derivation, independent review, computation and exposition.}}
\date{20 November 2026 --- research draft; priority not established}
\begin{document}\maketitle

\begin{abstract}
We test nonlinear distribution constructions and adaptive moment
correction against the Hadamard extraction target. Minimal binary
encodings of the known order-twelve and order-twenty matrices need
maximum Boolean degree exactly three, and binary-rank constraints
force growing degree for large orders congruent to four modulo eight.
We give a positivity-controlled support-weight correction and show
why no reweighting can repair a failed minimal-size support.
Separately, the four-bit pairwise-independent moment polytope has an
exact R\'enyi-order transition with competing strict local optima.
A concrete partial-matrix augmentation shows why odd-kernel shortness
is not a monotone construction potential. These are benchmarks and
conditional tools; no near-extremal distribution is constructed at
an unknown Hadamard order. The full conjecture remains open, and
novelty has not been established by a comprehensive literature review.
\end{abstract}
We retain the earlier target: for $N\ge8$ divisible by four, a law on
$N-4$ unbiased pairwise-independent sign bits with collision at least
$1/N-1/(1024N^4)$ yields a Hadamard matrix. Its proof is preserved in
\path{../round6/collision_stability.tex}; the exact residue-sensitive
entropy theorem is in \path{../round5/collision_entropy.tex}.
Neither characterization assumes attainment in every dimension.
\section{Minimal binary seeds require substantial nonlinear observables}
\label{sec:nonlinear-seeds}
We distinguish Boolean algebraic degree over $\mathbb F_2$ from real
Walsh degree. A sign observable is written $(-1)^{f(z)}$, and its degree
means the degree of the unique squarefree Boolean polynomial $f$.
The constructions below encode the classical Paley matrices; they do
not give new Hadamard orders. Paley's primary paper is
\emph{On Orthogonal Matrices} (1933),
\url{https://doi.org/10.1002/sapm1933121311}. The standard binary-code
viewpoint also appears in Orrick, \emph{Switching operations for Hadamard
matrices}, \url{https://arxiv.org/abs/math/0507515}.
All encoding bounds below are proved explicitly; novelty is unestablished.

\begin{proposition}[Affine observable budget]
\label{prop:nonlinear-affine-budget}
Suppose an order-$N$ normalized Hadamard matrix is encoded by distinct
labels $S\subset\mathbb F_2^r$, with uniform weight $1/N$ on each label.
If $N\equiv4\pmod8$ and $r<N-2$, the number $k$ of nonconstant columns
that extend to affine Boolean functions satisfies
\[
 k\le\min\{r,\,2^r-N-1\}\qquad(N<2^r).
\]
Thus a minimal four-bit encoding of $H_{12}$ needs at least eight
nonaffine observables, and a minimal five-bit encoding of $H_{20}$
needs at least fourteen.
\end{proposition}
\begin{proof}
Write $B=(J-H_{\rm core})/2$. Over $\mathbb F_2$, the core rows have
even weight $N/2$ and distinct rows intersect in odd weight $N/4$.
Therefore $BB^T=J+I$, of rank $N-2$. Also $B\mathbf1=0$, so
$\operatorname{rank}B=N-2$ and its only column relation is the sum
of all $N-1$ columns. Any proper subset is independent. Affine functions
vanishing at the label of the normalized first row form a space of
dimension $r$. Hence at most $r$ columns can be affine.

For the other bound, affine sign functions giving distinct nonconstant
orthogonal columns have distinct nonzero linear parts. Together with
one, they are orthogonal on the whole binary cube. Their inner products
vanish on $S$, so also on its complement $T$. Thus $k+1$ nonzero
orthogonal real functions live on $|T|=2^r-N$ points, giving the bound.
\end{proof}

Let $\mathrm{RM}(d,r)$ denote the space of Boolean polynomials of degree
at most $d$ on $r$ bits; its dimension is $\sum_{j=0}^d\binom rj$.
\begin{lemma}[Complement criterion for low-degree extensions]
\label{lem:nonlinear-extension}
Every Boolean function on $S\subset\mathbb F_2^r$ has a degree-at-most-$d$
extension if and only if there is no nonzero polynomial of degree at
most $r-d-1$ supported entirely on $S$. In particular, cubic extensions
exist on any proper subset of the four-bit cube. On five bits they
exist whenever the complement of $S$ affinely spans the cube.
\end{lemma}
\begin{proof}
The binary inner product of monomials of total combined degree less
than $r$ is zero: summing over any missing variable gives zero modulo
two. Dimension counting then gives
$\mathrm{RM}(d,r)^\perp=\mathrm{RM}(r-d-1,r)$.
A nonzero left-kernel vector of restriction to $S$, extended by zero
off $S$, is exactly such a dual polynomial. The particular cases use
that a nonzero constant cannot vanish on a nonempty complement, and
that an affine function vanishing on an affine spanning set is zero.
\end{proof}

\begin{theorem}[Sharp cubic benchmarks at orders twelve and twenty]
\label{thm:nonlinear-cubic}
For $N=12$ and $N=20$, the least possible maximum Boolean degree of
all $N-1$ observables, among encodings on the minimal number
$r=\lceil\log_2N\rceil$ of binary seed bits, is exactly three.
The affine budgets above are simultaneously sharp: the examples have
exactly three affine observables at order twelve and five at order twenty.
\end{theorem}
\begin{proof}
For $N\equiv4\pmod8$, the normalized binary columns together with the
constant function span a space of dimension $N-1$: the binary columns
have rank $N-2$ and vanish at the first row. Thus any common degree
bound $d$ requires
$\sum_{j=0}^d\binom rj\ge N-1$.
At $N=20,r=5$, quadratics have dimension $16<19$, proving the lower
bound three.

At $N=12,r=4$, the quadratic dimension is eleven, exactly the required
rank. Suppose all observables were quadratic. Restriction of
$\mathrm{RM}(2,4)$ to the twelve labels would have rank eleven and equal
the span of the binary Hadamard columns and the constant function.
Every vector in the latter span has even weight: the column weights
are six and the constant has weight twelve. Hence all coordinate
functions have even weight on $S$. Their full-cube weights are even,
so the four omitted labels have xor zero. Four distinct labels with xor
zero form an affine plane. Its indicator is a nonzero quadratic
polynomial vanishing on $S$, contradicting the restriction rank eleven.
This rules out quadratics for every labeling, without exhaustive search.

For explicit upper bounds, use Paley's normalized order-$p+1$ matrix
for $p=11,19$:
\[
 H_{\infty,j}=H_{i,\infty}=1,\qquad
 H_{i,j}=\chi(j-i)-\mathbf1_{i=j},
\]
where $\chi$ is the Legendre symbol with $\chi(0)=0$.
The first three nonconstant columns of $H_{12}$ give eight binary
patterns, four occurring twice and four once. Label rows by these
three bits and a fourth discriminator, choosing zero for the first
occurrence and one for the second. In row order $(\infty,0,\ldots,10)$
the resulting four-bit integer labels are
\[
 0,5,3,6,13,11,7,14,4,8,1,2.
\]
The first three observables are affine. Every other observable has a
cubic extension by the lemma. The affine budget proves three is maximal.

The first five nonconstant columns of $H_{20}$ separate all twenty
rows, with five-bit labels
\[
 0,13,27,22,12,25,19,7,15,30,29,26,21,10,20,8,16,1,3,6.
\]
The complement contains the affine basis
$2,4,5,9,11,17$: their differences from $2$ have binary rank five.
The lemma gives cubic extensions of every observable, retaining the
first five as coordinate functions. The affine budget proves optimality.
All displayed finite certificates and extension polynomials are checked
by original exact binary elimination in the accompanying script.
\end{proof}

\paragraph{A reusable nonlinear seed rule and its limit.}
For any odd prime power $q\equiv3\pmod4$, Paley's quadratic-character
formula over $\mathbb F_q$ supplies an extremal $(q+1)$-atom distribution.
Inject its row labels into a binary cube, use the uniform measure on
those labels, and solve the Boolean restriction equations for the
observables. Lemma~\ref{lem:nonlinear-extension} precisely certifies when
a chosen degree budget can express all outputs. The order-twelve and
order-twenty examples show that cubic outputs genuinely cross the affine
seed obstruction at minimal seed length; they are exact extremal laws,
not approximate moment fits. They still inherit Paley's restriction on
$q$, so this rule supplies no distribution at a currently unknown order.

\begin{corollary}[Fixed algebraic degree cannot suffice asymptotically]
\label{cor:nonlinear-growing-degree}
For any Hadamard order $N\ge12$ with $N\equiv4\pmod8$, and minimal
binary seed length $r=\lceil\log_2N\rceil$, every encoding has maximum
Boolean degree at least $\lceil r/2\rceil$. Conversely, every injection
of its rows into the $r$-bit cube admits extensions of degree at most
$r-1$. Thus a fixed-degree minimal-seed model cannot represent extremal
laws in arbitrarily large such orders, even where Hadamard matrices
are already known.
\end{corollary}
\begin{proof}
The rank bound requires $\sum_{j=0}^d\binom rj\ge N-1$.
Here $N-1>2^{r-1}$, while binomial symmetry gives
$\sum_{j=0}^{\lfloor(r-1)/2\rfloor}\binom rj\le2^{r-1}$.
This proves the lower bound. The proper subset $S$ has nonempty
complement, and Lemma~\ref{lem:nonlinear-extension} with $d=r-1$
gives the upper bound.
\end{proof}

\section{Adaptive support corrections and a support obstruction}

Let $S=\{x_1,\ldots,x_s\}\subseteq\{\pm1\}^d$ consist of distinct
vectors. Put $k=1+d+\binom d2$ and let $A$ be the $k\times s$ matrix
whose rows are evaluations on $S$ of the characters of degrees zero,
one and two. Set $b=(1,0,\ldots,0)^T$. A probability vector $p$ is
unbiased and pairwise independent precisely when $Ap=b$.

\begin{theorem}[Actual positive correction with an explicit margin]
Suppose $A$ has full row rank, and define
\[
 R=A^T(AA^T)^{-1},\quad
 \gamma=\max_j\sum_a|R_{ja}|,\quad
 m=Ap-b,\quad \tau=\|m\|_\infty.
\]
For any probability vector $p$ with $p_{\min}\ge\gamma\tau$, the
vector $q=p-Rm$ is a nonnegative probability vector on the same
support and satisfies $Aq=b$. In particular this is an exact
distribution correction, not an unconstrained projection being
mistaken for a probability law. It obeys
\[
 \|q-p\|_\infty\le\gamma\tau,\quad
 d_{\rm TV}(p,q)\le s\gamma\tau/2,\quad
 C(q)\ge C(p)-2\gamma\tau.
\]
\end{theorem}
\begin{proof}
Full row rank gives $AR=I_k$, so $Aq=b$. The definition of $\gamma$
gives $|(Rm)_j|\le\gamma\tau$, proving positivity and both distance
bounds. The degree-zero equation of $Aq=b$ says $\sum q_j=1$.
Finally write $h=q-p$; then
$C(q)=C(p)+2\sum p_jh_j+\sum h_j^2\ge C(p)-2\|h\|_\infty$.
\end{proof}
All inverse entries, rank conditions and margins are exact rational
calculations for rational input weights. If $A$ is not full row rank,
one may instead solve on its column span, but first must verify
$b\in\operatorname{range}A$; small residual alone does not imply this.
The positivity margin must still be checked. The theorem gives no
license to omit either condition.

For $d=n-1$, $n\equiv1\pmod4$, $N=n+3$, the extra hypotheses
\[
 C(p)\ge1/N,\qquad \gamma\tau\le1/(2048N^4)
\]
place $q$ in the round6 extraction window and thus allow exact
Hadamard extraction. This is a sufficient criterion on a supplied
distribution and support, not a construction satisfying it in every
unknown order. The conditioning of the support enters explicitly
through $\gamma$, and tiny reservoir weights enter through $p_{\min}$.

\begin{theorem}[No adaptive reweighting on a minimal-size support]
Let $N\ge8$ be divisible by four, $d=N-4$, and $S$ have at most $N$
distinct sign vectors. If an unbiased pairwise-independent law exists
on $S$, its positive support has exactly $N$ members and its weights
are all $1/N$. Its augmented sign rows complete to a Hadamard matrix
of order $N$.
Consequently, if the uniform law on a given $N$-point set is not
isotropic, no reweighting on that same set can make it isotropic.
\end{theorem}
\begin{proof}
If the positive support has size $t\le N$, Cauchy--Schwarz gives
$C\ge1/t\ge1/N$. The residue-sensitive collision bound proved in
round5 gives $C\le1/N$. Hence $t=N$ and equality in Cauchy--Schwarz
forces uniform weights. The equality case of that same bound gives
the rank-three tetrahedral completion to a Hadamard matrix of order
$N$.
\end{proof}
This obstruction applies to every adaptive weight correction on the
specified support, rather than only the scalar Fourier buffers of
round7. Such a correction must change or enlarge a failed minimal
support. Conversely the full-row-rank correction theorem necessarily
uses $s\ge k$, which exceeds $N$ when $d=N-4\ge4$. Its hypotheses
therefore address genuinely richer supports and cannot covertly
circumvent the minimal-support obstruction.
\section{An exact entropy-order transition on four bits}
\label{sec:renyi-phase}
An entropy-order continuation method would vary the objective in the
hope that the desired Hadamard distribution becomes easier to reach.
The smallest nontrivial moment polytope already has a complete,
explicit transition with persistent competing local optima. The
following analysis is self-contained; it does not establish novelty
or an all-order existence theorem.

Let $\mathcal P_4$ be the laws on four sign bits with zero means and
zero distinct pair correlations. For $0<\alpha<\infty$, $\alpha\ne1$,
define
\[
 H_\alpha(p)=\frac{\log\sum_xp_x^\alpha}{1-\alpha},
\]
with $H_1$ the Shannon entropy and all logarithms natural.
Define two vertex families: $\mathcal U$ consists of the ten uniform
laws obtained by fixing one triple parity or the quadruple parity;
$\mathcal W$ consists of the sixteen sign translations of
\[
 w_x=\begin{cases}
 1/6,&k(x)=0,\\
 1/12,&k(x)\in\{2,3\},\\
 0,&k(x)\in\{1,4\},
 \end{cases}
\]
where $k$ counts negative coordinates.

\begin{theorem}[Complete four-bit entropy diagram]\label{thm:renyi-phase}
There is a unique $\alpha_*\in(6,7)$ satisfying
\[
 (4/3)^{\alpha_*}+10(2/3)^{\alpha_*}=8.
\]
For every finite $\alpha>0$, the global minimizers of $H_\alpha$ on
$\mathcal P_4$ are exactly $\mathcal U$ when $\alpha<\alpha_*$,
exactly $\mathcal W$ when $\alpha>\alpha_*$, and their union at
$\alpha=\alpha_*$. Explicitly, the minimum is the smaller of
\[
 \log8,\qquad
 \frac{\log(2^\alpha+10)-\alpha\log12}{1-\alpha},
\]
with the second expression interpreted by continuity at one.
Every law in both families remains a strict local minimum for every
finite $\alpha>1$. Thus varying the entropy order does not force
the competing local minima to disappear at the global transition.
\end{theorem}

\begin{proof}
Fourier expansion writes every feasible law uniquely as
\[
 16p_x=1+\sum_{i=1}^4a_i\chi_{[4]\setminus\{i\}}(x)
                   +a_5\chi_{[4]}(x).
\]
The five displayed characters take exactly the sixteen sign patterns
$v\in\{\pm1\}^5$ with $\prod_i v_i=1$. Hence the parameter polytope is
\[
 P=\{a\in\mathbb R^5:1+v\cdot a\ge0\ \text{for all such }v\}.
\]
We first prove the exact vertex description
\[
 P=\operatorname{conv}\left(
 \{\pm e_i:1\le i\le5\}
 \cup\{s/3:s\in\{\pm1\}^5,\ \prod_i s_i=1\}\right).
\]
All listed points satisfy the inequalities. Conversely put
$b_i=|a_i|$ and $m=\min_i b_i$. If some coordinate is zero or the
product of coordinate signs is negative, an admissible $v$ gives
$v\cdot a=-\sum_i b_i$, so $\sum_i b_i\le1$ and $a$ lies in the
cross-polytope generated by the signed coordinate vectors.
Otherwise let $s_i=\operatorname{sign}(a_i)$, with $\prod s_i=1$.
The least admissible dot product is $-\sum_i b_i+2m$, so
$\sum_i b_i-2m\le1$. The representation
\[
 a=(3m)(s/3)+\sum_i(b_i-m)s_i e_i
\]
has nonnegative coefficients of total $\sum_i b_i-2m\le1$;
assign any remaining mass to zero, which is in the cross-polytope.
This proves the hull identity. The signed coordinate vectors are
unique extrema of coordinate functionals. Each $s/3$ uniquely maximizes
$s\cdot a$ among the listed points, attaining $5/3$ rather than at
most one elsewhere. Thus all twenty-six points are vertices.
They yield precisely $\mathcal U$ and $\mathcal W$ above.

For $\alpha>1$, minimizing entropy means maximizing the strictly
convex function $\sum p^\alpha$; for $0<\alpha<1$ it means minimizing
that strictly concave function. At one Shannon entropy is strictly
concave. In every case a nonvertex cannot minimize, so it suffices
to compare the two weight multisets. For $\alpha\ne1$ their power sums
differ in sign according to
\[
 \phi(\alpha)=(4/3)^\alpha+10(2/3)^\alpha-8.
\]
This function is strictly convex. It has a root at one, and exact
rational evaluation gives $\phi(6)<0<\phi(7)$. There is exactly
one other root, $\alpha_*\in(6,7)$; the signs are positive before
one, negative between one and $\alpha_*$, and positive afterwards.
Accounting for the sign of $1-\alpha$ proves the stated comparison.
At one the weighted entropy is
$\log12-(\log2)/6>\log8$. Strict convexity or concavity also excludes
nonvertex mixtures at the crossing.

It remains to verify local persistence, not merely stationarity.
Fix finite $\alpha>1$. For a weighted vertex put
$u=(1/6)^{\alpha-1}$ and $b=(1/12)^{\alpha-1}$. The degree-two function
\[
 q(k)=b+\frac{u-b}{6}(k-2)(k-3)
\]
equals $w_x^{\alpha-1}$ on its positive support and is strictly
positive elsewhere. Its expectation is fixed by the moment constraints.
It therefore strictly exposes the gradient at $w$, unless a feasible
law has the same support. There is no other feasible law on that
support: a second one would allow a small continuation from $w$ in
the opposite direction and contradict that $w$ is a vertex.
For a uniform vertex the same argument uses the constant polynomial
equal to its nonzero gradient value.

For clarity, strict linear exposure suffices here because the feasible
set is a finite polytope. The negative gradient slope divided by the
distance to each other vertex has a positive minimum. Convex
decomposition and the triangle inequality give a uniform negative
linear bound in every feasible direction. Differentiability of
$\sum p^\alpha$ for $\alpha>1$ makes its remainder $o(\|p-w\|)$,
proving strict local maximality of the power sum and hence strict
local minimality of entropy. This argument includes $1<\alpha<2$,
without requiring a bounded Hessian at zero coordinates.
\end{proof}

The order-one root of $\phi$ reflects normalization and is not an
entropy transition. The genuine transition lies between six and seven.
As $\alpha$ tends to infinity the weighted optimum approaches
$\log6$, whereas the collision and Shannon optima are $\log8$.
The local-persistence statement explicitly excludes the endpoint
$\alpha=\infty$.

This example is a calibrated obstruction to using entropy-order
continuation alone as a global attainment argument. It does not
exclude nonlocal moves or an independently justified global method.
It also illustrates a useful geometric reduction: the four-bit
moment polytope is the polar of a five-dimensional parity half-cube,
and its entire optimization diagram follows from its explicit hull.
\section{A concrete repair and a failure of odd-coset monotonicity}
\label{sec:odd-coset-repair}

We continue the $12\times24$ terminal example whose all-odd integer
kernel minimum was certified to be $40$, rather than the sign-vector
value $24$. Exact search finds a one-row exchange, substantially more
efficient than rounding the displayed short odd vector. It also shows
that the odd-coset minimum is not a monotone potential for augmentation:
an explicitly larger frame has no all-odd integer kernel vector at all.
The broader distinction between arithmetic and sign completion is
classical; see Verheiden, \emph{Integral and rational completions of
Hadamard matrices},
\url{https://doi.org/10.1016/0097-3165(80)90028-X}.
The finite certificates and deductions below are independently checked.

Use zero-based coordinates and encode a sign row by the integer having
bit $j$ equal to one exactly when coordinate $j$ is $-1$. Let $P$
have rows with masks
\[
\begin{gathered}
0,16773120,16519104,9656012,13838682,14783914,\\
1944718,11409518,15230494,5921970,3487346,13018356.
\end{gathered}
\]
The rows are orthogonal. The previously studied odd kernel vector is
\[
\begin{split}
v={}&(-1,-1,3,-1,1,1,-3,-1,1,1,-1,1,\\
&1,-1,-1,-1,1,-1,1,-1,1,-1,1,1),
\end{split}
\]
with $Pv=0$ and $\|v\|^2=40$. The two exceptional coordinates are
$2,6$. Its coordinatewise sign rounding is orthogonal to only seven
of the twelve original rows; it violates precisely rows
$\{2,4,6,8,10\}$. Keeping this particular rounded candidate would
therefore require discarding five original rows.

\begin{proposition}[Exact one-row augmentation]
The minimum number of original rows that must be deleted from $P$
to permit an augmenting exchange is exactly one. Specifically delete
row index $1$, with mask $16773120$, and insert rows with masks
\[
3198332,\qquad6782212.
\]
The resulting thirteen rows are orthogonal. More strongly, after this
same deletion the maximum number of mutually orthogonal inserted
rows is exactly five, attained by
\[
3198332,6782212,8735670,11139280,13957672.
\]
The resulting sixteen-row matrix $Q$ is terminal.
\end{proposition}
\begin{proof}
There is no sign row orthogonal to all twelve original rows, so zero
deletions cannot augment. This is certified by exact rational kernel
elimination and enumeration of its $2^{11}=2048$ information-coordinate
sign assignments, up to negation.
Deleting row $1$ gives exactly sixteen antipodal sign-candidate classes.
The checker lists all their masks. Two masks are adjacent exactly when
their Hamming distance is $12$. The five displayed masks form a clique,
and exhaustive recursion on this sixteen-vertex graph finds no
six-clique. Thus five is the exact maximum. All row norms and pairwise
inner products are checked as integers.
If another sign row extended $Q$, it would belong to the same candidate
pool and enlarge the five-clique, which is impossible.
\end{proof}

Deleting row $0$, in contrast, gives only two candidate classes and
no orthogonal pair. Both deletion choices close the odd-coset norm gap
to $24$, since the deleted original sign row itself enters the kernel.
Therefore closing that gap cannot distinguish successful augmentation
from a deletion permitting only replacement.

\begin{proposition}[Augmentation can destroy odd integer feasibility]
The larger frame $Q$ has no all-odd integer kernel vector. In particular
the minimum squared norm in the odd integer coset increases from
$40$ for $P$ to $+\infty$ for $Q$, under an explicit row-augmenting
exchange.
\end{proposition}
\begin{proof}
Let $f=e_0+e_9-e_{11}$. Exact calculation gives
\[
 24\|f\|^2-\sum_{q\text{ row of }Q}(q\cdot f)^2=0.
\]
For orthogonal rows of norm $\sqrt{24}$ this means that $f$ lies in
their rational row span. Its coordinate sum is odd. Any all-odd integer
vector $z$ has $f\cdot z$ odd, whereas $Qz=0$ would force
$f\cdot z=0$. This contradiction proves the assertion and also gives
a small direct certificate of terminality for $Q$.
\end{proof}

\subsection{What coordinatewise rounding does guarantee}
\begin{proposition}[Rounding defect]
Let $P$ have $r$ orthogonal sign rows of length $N$, and suppose
$v\in\ker P$ has entries in $\{\pm1,\pm3\}$. Let $w$ be its
coordinatewise sign, and let $t$ count the coordinates of magnitude
$3$. Define $e_j=w_j$ on those coordinates and $e_j=0$ otherwise.
Then $v=w+2e$ and $Pw=-2Pe$.
If $t$ is odd, $w$ is orthogonal to \emph{none} of the original rows.
If $t$ is even, the number $b$ of rows not orthogonal to $w$ satisfies
\[
 b\le\min\{r,Nt/4\}.
\]
Deleting exactly those $b$ rows permits inserting $w$ as a single
sign row. This need not be an augmenting exchange.
\end{proposition}
\begin{proof}
For every sign row $p$, $p\cdot e\equiv t\pmod2$. If $t$ is odd,
that inner product cannot vanish. If $t$ is even, each nonzero value
has magnitude at least $2$, while Bessel gives
$\sum_p(p\cdot e)^2\le N\|e\|^2=Nt$. Hence $4b\le Nt$.
The identity $Pw=-2Pe$ supplies all remaining claims.
\end{proof}

Thus a short odd vector supplies a deterministic candidate and a
computable deletion set, but its number of exceptional coordinates
alone has not here yielded a bounded augmenting-exchange theorem.
The explicit example has rounding deletion count five and actual
augmentation deletion count one. The earlier sixteen-row order-$24$
example has a certified three-row augmentation barrier and already
fails odd integer feasibility. These examples separate three issues:
parity feasibility, shortness of an odd kernel vector, and mutual
compatibility of enough new sign rows.

Finally, a global column sign change cannot repair any of these
obstructions. Multiplication by a diagonal sign matrix gives a
norm-preserving bijection on sign extensions and on odd integer kernel
vectors, and preserves every row inner product. Column permutations
have the same property. Such operations can simplify a search's
normalization, but cannot change its extension or exchange costs.

\section{The next existence step}
The nonlinear encodings construct exact extremal laws at classical
known orders and prove specific lower bounds on model complexity.
They do not supply an encoding whose uniform support is a Hadamard
design in an unknown order. The adaptive correction works only with
its explicit rank, range and positivity conditions; no input meeting
the extraction window at all orders has been proved.

The entropy transition and lattice repair give calibrated warnings.
Changing entropy order need not destabilize a nonglobal local
minimum, and a successful augmentation can increase odd-kernel
shortness to infinity. A global construction argument must justify
its choice of nonlinear supports, nonlocal moves, or a different
potential rather than assume those diagnostics converge.
All universal statements have written proofs; finite certificates are
checked by exact arithmetic. No Lean code or Git commit was made.
\end{document}
